% ============================================================================
% CAPÍTULO 3: FERRAMENTAS MODERNAS DE DESENVOLVIMENTO E IA
% Bioinformática para Biologia de Sistemas
% ============================================================================

% ============================================================================
% TEMPLATE BASE PARA APRESENTAÇÕES EM BEAMER
% Bioinformática para Biologia de Sistemas
% ============================================================================

\documentclass[12pt, aspectratio=169]{beamer}

% ============================================================================
% PACOTES ESSENCIAIS
% ============================================================================
\usepackage[utf8]{inputenc}
\usepackage[T1]{fontenc}
\usepackage[brazilian]{babel}
\usepackage{amsmath, amssymb, amsthm}
\usepackage{graphicx}
\usepackage{xcolor}
\usepackage{tikz}
\usetikzlibrary{arrows.meta, shapes, positioning, calc, decorations.pathreplacing, decorations.shapes, decorations.pathmorphing}
\usepackage{listings}
\usepackage{booktabs}
\usepackage{multirow}
\usepackage{hyperref}
\usepackage{chemfig}
\usepackage{chemarr}

% ============================================================================
% CONFIGURAÇÃO DO TEMA
% ============================================================================
\usetheme{Madrid}
\usecolortheme{default}

% Cores personalizadas
\definecolor{azulescuro}{RGB}{0, 51, 102}
\definecolor{azulclaro}{RGB}{51, 153, 255}
\definecolor{verdeacido}{RGB}{102, 204, 0}
\definecolor{laranjacel}{RGB}{255, 153, 51}
\definecolor{roxodna}{RGB}{153, 51, 255}

\setbeamercolor{structure}{fg=azulescuro}
\setbeamercolor{palette primary}{bg=azulescuro,fg=white}
\setbeamercolor{palette secondary}{bg=azulclaro,fg=white}
\setbeamercolor{palette tertiary}{bg=azulescuro,fg=white}
\setbeamercolor{palette quaternary}{bg=azulclaro,fg=white}
\setbeamercolor{block title}{bg=azulescuro,fg=white}
\setbeamercolor{block body}{bg=azulclaro!10,fg=black}
\setbeamercolor{block title alerted}{bg=red!80,fg=white}
\setbeamercolor{block body alerted}{bg=red!10,fg=black}
\setbeamercolor{block title example}{bg=verdeacido!80,fg=white}
\setbeamercolor{block body example}{bg=verdeacido!10,fg=black}

% ============================================================================
% CONFIGURAÇÃO DE CÓDIGO
% ============================================================================
\lstset{
    language=Python,
    basicstyle=\ttfamily\small,
    keywordstyle=\color{azulescuro}\bfseries,
    commentstyle=\color{verdeacido},
    stringstyle=\color{laranjacel},
    numbers=left,
    numberstyle=\tiny\color{gray},
    stepnumber=1,
    numbersep=5pt,
    backgroundcolor=\color{white},
    showspaces=false,
    showstringspaces=false,
    showtabs=false,
    frame=single,
    rulecolor=\color{black},
    tabsize=2,
    captionpos=b,
    breaklines=true,
    breakatwhitespace=false,
    escapeinside={\%*}{*)},
    xleftmargin=10pt,
    xrightmargin=10pt
}

% Definir estilo para R
\lstdefinestyle{Rstyle}{
    language=R,
    basicstyle=\ttfamily\small,
    keywordstyle=\color{azulescuro}\bfseries,
    commentstyle=\color{verdeacido},
    stringstyle=\color{laranjacel}
}

% Definir estilo para MATLAB
\lstdefinestyle{MATLABstyle}{
    language=Matlab,
    basicstyle=\ttfamily\small,
    keywordstyle=\color{azulescuro}\bfseries,
    commentstyle=\color{verdeacido},
    stringstyle=\color{laranjacel}
}

% ============================================================================
% COMANDOS PERSONALIZADOS
% ============================================================================

% Comando para equações destacadas
\newcommand{\destaque}[1]{%
    \begin{alertblock}{}
        \begin{center}
            #1
        \end{center}
    \end{alertblock}
}

% Comando para definições
\newcommand{\definicao}[2]{%
    \begin{block}{Definição: #1}
        #2
    \end{block}
}

% Comando para exemplos
\newcommand{\exemplo}[2]{%
    \begin{exampleblock}{Exemplo: #1}
        #2
    \end{exampleblock}
}

% Comando para observações importantes
\newcommand{\importante}[1]{%
    \begin{alertblock}{Importante!}
        #1
    \end{alertblock}
}

% Comandos matemáticos úteis
\newcommand{\R}{\mathbb{R}}
\newcommand{\N}{\mathbb{N}}
\newcommand{\Z}{\mathbb{Z}}
\newcommand{\dif}{\mathrm{d}}
\newcommand{\parcial}[2]{\frac{\partial #1}{\partial #2}}
\newcommand{\derivada}[2]{\frac{\dif #1}{\dif #2}}
\newcommand{\vect}[1]{\boldsymbol{#1}}

% Comandos biológicos
\newcommand{\gene}[1]{\textit{#1}}
\newcommand{\proteina}[1]{\textsc{#1}}
\newcommand{\especie}[1]{\textit{#1}}

% ============================================================================
% CONFIGURAÇÃO DE RODAPÉ
% ============================================================================
\setbeamertemplate{footline}{
  \leavevmode%
  \hbox{%
  \begin{beamercolorbox}[wd=.33\paperwidth,ht=2.25ex,dp=1ex,center]{author in head/foot}%
    \usebeamerfont{author in head/foot}\insertshortauthor
  \end{beamercolorbox}%
  \begin{beamercolorbox}[wd=.34\paperwidth,ht=2.25ex,dp=1ex,center]{title in head/foot}%
    \usebeamerfont{title in head/foot}\insertshorttitle
  \end{beamercolorbox}%
  \begin{beamercolorbox}[wd=.33\paperwidth,ht=2.25ex,dp=1ex,right]{date in head/foot}%
    \usebeamerfont{date in head/foot}\insertshortdate{}\hspace*{2em}
    \insertframenumber{} / \inserttotalframenumber\hspace*{2ex}
  \end{beamercolorbox}}%
  \vskip0pt%
}

% ============================================================================
% CONFIGURAÇÃO DE NAVEGAÇÃO
% ============================================================================
\setbeamertemplate{navigation symbols}{}

% ============================================================================
% CONFIGURAÇÃO DE ITEMIZE/ENUMERATE
% ============================================================================
\setbeamertemplate{itemize items}[circle]
\setbeamertemplate{enumerate items}[default]

% ============================================================================
% FIM DO TEMPLATE
% ============================================================================


\title[Cap. 3: IDEs e IA]{Capítulo 3: Ferramentas Modernas de Desenvolvimento e IA}
\subtitle{O Fim do "Terminal Assustador" e a Revolução dos Copilotos}
\author{Ney Lemke}
\institute{DBF-Unesp}
\date{\today}

\begin{document}

% ============================================================================
% SLIDE DE TÍTULO
% ============================================================================
\begin{frame}
\titlepage
\end{frame}

% ============================================================================
% SUMÁRIO
% ============================================================================
\begin{frame}{Sumário}
\tableofcontents
\end{frame}

% ============================================================================
% INTRODUÇÃO
% ============================================================================
\section{Introdução}

\begin{frame}{O fim do ``Terminal Assustador''}
\begin{itemize}
    \item Para não-programadores, aprender sintaxe é muitas vezes a barreira mais frustrante.
    \item O paradigma antigo: "Decorar comandos e lutar contra erros de sintaxe".
    \item \textbf{O novo paradigma:} Ferramentas e Inteligência Artificial integradas ao ambiente de trabalho.
\end{itemize}

\vspace{0.5cm}
\begin{block}{O Piloto e o Copiloto}
A máquina (Copiloto) lida com a \textbf{sintaxe, automação e sugestões}. Você, o pesquisador (Piloto), foca na \textbf{lógica biológica e interpretação científica dos dados}.
\end{block}
\end{frame}

% ============================================================================
% AMBIENTE DE DESENVOLVIMENTO
% ============================================================================
\section{O Ambiente de Desenvolvimento (IDE)}

\begin{frame}{Ambiente de Desenvolvimento Integrado (IDE)}
\begin{columns}
\column{0.55\textwidth}
\textbf{O que é uma IDE?} \\
Um software que centraliza todas as ferramentas que um programador precisa:
\begin{itemize}
    \item Editor de texto com realce de sintaxe
    \item Terminal embutido
    \item Gerenciador de arquivos
    \item Depurador (Debugger)
\end{itemize}

\vspace{0.3cm}
\textbf{Visual Studio Code (VS Code)}
\begin{itemize}
    \item O padrão da indústria hoje.
    \item Altamente customizável via extensões (Python, R, Jupyter).
\end{itemize}

\column{0.4\textwidth}
\begin{exampleblock}{Dica Prática}
Nunca use o "Bloco de Notas". Uma IDE ajuda a capturar erros antes mesmo de você rodar o código.
\end{exampleblock}
\end{columns}
\end{frame}

\begin{frame}{Cursor e Forks Modernos}
\begin{itemize}
    \item \textbf{Cursor} é um \textit{fork} (uma modificação direta) do VS Code.
    \item Diferença: Ele foi construído com a \textbf{Inteligência Artificial no seu núcleo}.
    \item Você não apenas digita código, você "conversa" com o seu projeto.
\end{itemize}

\vspace{0.5cm}
\begin{block}{Por que usar o Cursor ou similares?}
Eles entendem não apenas o arquivo aberto, mas a base de dados inteira. Você pode perguntar: \textit{"Onde está o script que gera os gráficos de PCA?"}
\end{block}
\end{frame}

% ============================================================================
% REVOLUÇÃO DA IA
% ============================================================================
\section{A Revolução da IA no Código}

\begin{frame}{Modelos de Linguagem para Código}
\begin{columns}
\column{0.6\textwidth}
O que mudou nos últimos anos? O treinamento massivo de \textbf{LLMs (Large Language Models)} em bilhões de linhas de código público (GitHub, StackOverflow).

\vspace{0.3cm}
\textbf{Modelos líderes no desenvolvimento:}
\begin{itemize}
    \item \textbf{Claude 3.5 Sonnet} (Anthropic) - Atual estado da arte para codificação.
    \item \textbf{GPT-4o} (OpenAI) - Excelente lógica analítica.
    \item \textbf{Codex} - O modelo que originou as ferramentas de autocompletar.
\end{itemize}

\column{0.35\textwidth}
\begin{center}
A programação natural se tornou uma linguagem de compilação. O inglês/português agora também gera software.
\end{center}
\end{columns}
\end{frame}

\begin{frame}{Níveis de Assistência: Autocompletar vs Agentes}
\textbf{Nível 1: Autocompletar (Ex: GitHub Copilot)}
\begin{itemize}
    \item Você começa a digitar e ele prevê o restante da linha ou da função.
    \item \textit{Pró:} Muito rápido, economiza tempo de digitação repetitiva.
\end{itemize}

\vspace{0.3cm}
\textbf{Nível 2: Chat Integrado}
\begin{itemize}
    \item Uma janela lateral onde você cola um erro e pede para explicar.
\end{itemize}

\vspace{0.3cm}
\textbf{Nível 3: Agentes Autônomos (Ex: Antigravity, Claude Code)}
\begin{itemize}
    \item Ferramentas que lêem arquivos, tomam decisões, rodam testes e modificam o código sozinhas a partir de uma requisição ("Arrume esse bug de memória").
\end{itemize}
\end{frame}

% ============================================================================
% CASOS DE USO EM BIOINFORMÁTICA
% ============================================================================
\section{Casos de Uso na Bioinformática}

\begin{frame}[fragile]{Caso 1: Explicando Código Legado}
É comum herdar scripts sem documentação do laboratório.

\vspace{0.2cm}
\textbf{Sua requisição para a IA:}
\begin{quote}
"O que este script R faz exatamente? Ele normaliza os dados de RNA-seq usando qual método?"
\end{quote}

\vspace{0.2cm}
A IA pode destrinchar o script e adicionar comentários em português antes mesmo de você rodar.
\end{frame}

\begin{frame}[fragile]{Caso 2: Escrevendo Scripts Rápidos}
\textbf{Sua requisição para a IA:}
\begin{quote}
"Escreva um script em Python usando Biopython que leia `sequencias.fasta` e salve em `filtrados.fasta` apenas as sequências com mais de 500 pares de base."
\end{quote}

\vspace{0.2cm}
\begin{lstlisting}[language=Python, caption=Código gerado pela IA, basicstyle=\ttfamily\footnotesize]
from Bio import SeqIO

records = []
for record in SeqIO.parse("sequencias.fasta", "fasta"):
    if len(record.seq) > 500:
        records.append(record)

SeqIO.write(records, "filtrados.fasta", "fasta")
\end{lstlisting}
\end{frame}

% ============================================================================
% BOAS PRÁTICAS E PERIGOS
% ============================================================================
\section{Boas Práticas e Perigos}

\begin{frame}{Nem tudo são flores: Os Perigos da IA}
\begin{columns}
\column{0.5\textwidth}
\begin{alertblock}{Alucinações}
A IA tem a "necessidade" de responder. Se ela não sabe, pode \textbf{inventar} um pacote de R ou um teste estatístico que simplesmente não existe.
\end{alertblock}

\vspace{0.2cm}
\begin{block}{Privacidade de Dados}
Nunca cole sequências de genomas humanos não-anonimizados ou dados clínicos sensíveis de pacientes em prompts de IAs públicas.
\end{block}

\column{0.45\textwidth}
\begin{exampleblock}{O Conhecimento Fundacional}
Por que aprender o básico de lógica de programação?
\\ \vspace{0.2cm}
\textbf{Para auditar a IA.} Se você não sabe o que o código faz em essência, você não saberá identificar quando o resultado estiver cientificamente incorreto.
\end{exampleblock}
\end{columns}
\end{frame}

% ============================================================================
% CONCLUSÃO
% ============================================================================
\begin{frame}{Resumo}
\begin{itemize}
    \item IDEs como VS Code ou Cursor são indispensáveis para pesquisadores modernos.
    \item A Inteligência Artificial acelerou drasticamente a curva de aprendizado para biólogos computacionais.
    \item Utilize assistentes de código e agentes para ganhar velocidade e quebrar bloqueios lógicos.
    \item \textbf{Mas lembre-se:} A responsabilidade pelo rigor metodológico e biológico da pesquisa sempre será sua. A máquina escreve, mas quem assina o paper é você.
\end{itemize}
\end{frame}

\end{document}
